% Integration requirements: amsmath, amssymb, amsthm; theorem, lemma,
% proposition, corollary, remark, and proof environments.  No unprefixed
% macros are introduced.  Bibliography entries are in mixed_references.tex.
\providecommand{\mxLi}{\operatorname{Li}}

\section{Arbitrary depth with mixed integer inner orders}
\label{mx:section}

Question~11 of the preceding directional report asks for a finite
nested-polylogarithm reduction when the inner integer orders have both
signs.  Theorem~\ref{mx:main} supplies such a reduction while keeping the
outer order complex.  If exactly $q$ inner orders are positive, its
output has depth at most $q+1$, regardless of how many nonpositive
orders occur.  Its coefficients are finite rational operations,
root-of-unity filters, and generalized-harmonic polynomials.  Repeated
poles are included directly, and coincident colors are evaluated by an
explicit polynomial summation rule.

The shuffle algebra of iterated integrals and the reduction of rational
functions times hyperlogarithms are established tools; see
Panzer~\cite[Sections~2.1--2.3]{mx:Panzer}.  Colored depth-two Tornheim
reductions are also classical~\cite{mx:Zhao}.  The contribution here is
the explicit convolution formula with unrestricted outer order, its
depth bound, its coincident-color reduction rule, and the accompanying
meromorphic and normal-derivative prescription.  No claim of literature
priority for the underlying shuffle or rational-function algorithms is
made.

\subsection{Definitions and finite input data}

Let $s_j\in\mathbb Z$, $p_j\in\mathbb Z_{>0}$, and
$0<|X_j|\leq1$.  Define
\begin{equation}\label{mx:Z}
 Z_{\mathbf s,\mathbf p}(C;\mathbf X)
 =\sum_{m_1,\ldots,m_d\geq1}
  \frac{\prod_{j=1}^dX_j^{m_j}m_j^{-s_j}}
       {(p_1m_1+\cdots+p_dm_d)^C}.
\end{equation}
The half-plane
$\Re C>d+\sum_j\max(0,-s_j)$ is sufficient for absolute
convergence.  Every power of a positive integer uses its real
logarithm.  Put
\[
 I_+=\{j:s_j>0\},\qquad I_-=\{j:s_j\leq0\},\qquad
 q=|I_+|,
\]
and assume for the main theorem that both sets are nonempty.  Factor
the ordinary coefficient generating function as
\begin{equation}\label{mx:factor}
 \sum_{n\geq1}r_nt^n=R(t)P(t),\qquad
 R(t)=\prod_{j\in I_-}\mxLi_{s_j}(X_jt^{p_j}),\quad
 P(t)=\prod_{j\in I_+}\mxLi_{s_j}(X_jt^{p_j}).
\end{equation}
Then $Z(C)=\sum_{n\geq1}r_nn^{-C}$ in the stated half-plane.

For each $j\in I_-$ write $a_j=-s_j$.  The identity
\[
 \mxLi_{-a}(z)=\left(z\frac{d}{dz}\right)^a\frac z{1-z}
\]
shows that $R$ is rational and bounded at infinity.  Its finite partial
fractions have the form
\begin{equation}\label{mx:partial}
 R(t)=c_\infty+
       \sum_\alpha\sum_{h=1}^{M_\alpha}
          \frac{c_{\alpha,h}}{(1-\alpha t)^h},\qquad
 M_\alpha=\sum_{\substack{j\in I_-\\\alpha^{p_j}=X_j}}(1-s_j).
\end{equation}
The index $\alpha$ runs over the distinct roots of the indicated
equations.  Pole cancellation may make some coefficients zero.  Even
at repeated roots the coefficient formula is simply
\begin{equation}\label{mx:polecoeff}
 c_{\alpha,h}
 =[u^{M_\alpha-h}]\,u^{M_\alpha}R((1-u)/\alpha).
\end{equation}
The function inside the coefficient extraction is holomorphic at zero.
Because $I_-$ is nonempty, $R(0)=0$, so
\begin{equation}\label{mx:constantcancel}
 c_\infty+\sum_{\alpha,h}c_{\alpha,h}=0.
\end{equation}
Finally set
\begin{equation}\label{mx:e}
 e_{h,r}=[v^r]\prod_{\nu=1}^{h-1}(1+v/\nu)
 \quad(0\leq r<h).
\end{equation}
Thus $e_{h,0}=1$, $e_{h,1}=H_{h-1}$, and
$e_{h,2}=\{H_{h-1}^2-H_{h-1}^{(2)}\}/2$.  In every degree the
coefficients follow from
\[
 \prod_{\nu=1}^{h-1}(1+v/\nu)
 =\exp\!\left(\sum_{k\geq1}
          \frac{(-1)^{k+1}}kH_{h-1}^{(k)}v^k\right).
\]

We use the descending-index convention
\begin{equation}\label{mx:mpl}
 \mxLi_{u_1,\ldots,u_k}(z_1,\ldots,z_k)
 =\sum_{n_1>\cdots>n_k\geq1}
       \frac{z_1^{n_1}\cdots z_k^{n_k}}
            {n_1^{u_1}\cdots n_k^{u_k}}.
\end{equation}
The first order will remain variable.  Intermediate trailing orders
may be nonpositive integers; Lemma~\ref{mx:eliminate} removes all such
orders.  For root choices $\mathbf a=(a_j)_{j\in I_+}$ with
$a_j^{p_j}=X_j$, form the words
\[
 0^{s_j-1}a_j\qquad(j\in I_+).
\]
Their shuffle, with each original position retained as a labeled
position before equal words are collected, is a finite formal sum
$\sum_w\mu_w w$.  Every resulting word has exactly $q$ nonzero
letters and ends in a nonzero letter; hence it has a unique form
\begin{equation}\label{mx:word}
 w=0^{k_1-1}b_1\,0^{k_2-1}b_2\cdots0^{k_q-1}b_q,
 \qquad k_i\geq1.
\end{equation}
The sequence $(b_1,\ldots,b_q)$ lists the selected root colors in the
order in which they occur in that shuffle.  Repeated colors do not
alter this rule or suppress their shuffle multiplicities.

\subsection{The mixed-sign reduction theorem}

\begin{theorem}[Finite mixed-sign reduction and every normal jet]
\label{mx:main}
Let the preceding data be fixed, with $q\geq1$ and $I_-\ne\varnothing$.
Put $K=\prod_{j\in I_+}p_j^{s_j-1}$.  Then
\begin{align}\label{mx:mainformula}
 Z(C)={}&K\sum_{\substack{\mathbf a\\a_j^{p_j}=X_j}}
       \sum_w\mu_w\sum_\alpha\sum_{h=1}^{M_\alpha}
           c_{\alpha,h}
       \sum_{r=0}^{h-1}e_{h,r}\sum_{\ell=0}^r
           (-1)^\ell\binom r\ell\notag\\[-1mm]
 &\hspace{10mm}\cdot
 \mxLi_{C-r+\ell,\,k_1-\ell,\,k_2,\ldots,k_q}
 \left(\alpha,\frac{b_1}\alpha,
                  \frac{b_2}{b_1},\ldots,\frac{b_q}{b_{q-1}}\right).
\end{align}
This is an identity of meromorphic functions of $C$ on the whole
complex plane.  It reduces, by finite algebra, to a sum of terms
\begin{equation}\label{mx:positiveoutput}
 A\,\mxLi_{C+v,\,t_2,\ldots,t_k}(z_1,\ldots,z_k),
 \qquad v\in\mathbb Z,\quad t_i\in\mathbb Z_{>0},\quad k\leq q+1.
\end{equation}
All coefficients and colors are independent of $C$.  Consequently
every derivative $\partial_C^N Z$ is obtained by differentiating only
the first polylogarithm order in either finite formula.  If no $X_j$
equals one, the continuation is entire and these derivatives are
ordinary entire functions.  When poles occur, the identities and
derivatives are meromorphic identities with the pole prescription in
Theorem~\ref{mx:poles} below.
\end{theorem}

\begin{proof}
The root filter, valid first as an identity of power series, is
\begin{equation}\label{mx:rootfilter}
 \mxLi_s(Xt^p)=p^{s-1}\sum_{a^p=X}\mxLi_s(at).
\end{equation}
Indeed the sum of the $n$th powers of those roots vanishes unless
$p\mid n$, and equals $pX^{n/p}$ otherwise.  Thus no choice of a
particular $p$th-root branch enters the result.

For completeness, take the iterated-integral forms
$\omega_0=dt/t$ and $\omega_a=a\,dt/(1-at)$, and let
$\mathcal H_w$ denote the corresponding integral from zero with
outermost letter first.  All words used here end in a nonzero letter,
so these integrals require no tangential-base-point regularization.
Expansion of the kernels into geometric series proves
\begin{equation}\label{mx:hyperlog}
 \mathcal H_w(t)=
 \mxLi_{k_1,\ldots,k_q}
       (b_1t,b_2/b_1,\ldots,b_q/b_{q-1}).
\end{equation}
Partitioning the product of integration simplices by all total
orderings of their integration variables proves the shuffle product,
including its multiplicities when letters coincide.  Applying this
to \eqref{mx:rootfilter} gives
\[
 P(t)=K\sum_{\mathbf a}\sum_w\mu_w\mathcal H_w(t).
\]

It remains to compute one rational-factor convolution.  Write
$H(t)=\mathcal H_w(t)=\sum_{m\geq1}h_mt^m$.  For $n\geq m$,
\begin{align}\label{mx:binomial}
 [t^{n-m}](1-\alpha t)^{-h}
 &=\alpha^{n-m}\binom{n-m+h-1}{h-1}\notag\\
 &=\alpha^{n-m}\sum_{r=0}^{h-1}e_{h,r}
                   \sum_{\ell=0}^r(-1)^\ell\binom r\ell
                            n^{r-\ell}m^\ell.
\end{align}
The diagonal $n=m$ contributes $h_n$ exactly once.  The strict part
$n>m$, followed by the nested indices already in $h_m$, is precisely
the multiple polylogarithm in \eqref{mx:mainformula}.  The diagonal
contributions from every rational partial fraction and from
$c_\infty H$ have combined coefficient
$c_\infty+\sum_{\alpha,h}c_{\alpha,h}=R(0)=0$.  Thus they cancel
exactly, leaving the displayed formula.

This reasoning proves an identity of coefficient sequences.  All
series converge locally normally in $C$ if $|X_j|<1$: the cumulative
colors in every term of \eqref{mx:mainformula} are
$(\alpha,b_1,\ldots,b_q)$, all strictly inside the disk, so the
coefficients have exponential decay times a polynomial.  For boundary
colors the same coefficient identity can first be summed in a
sufficiently far right half-plane.  Lemma~\ref{mx:eliminate} produces
\eqref{mx:positiveoutput}; each elimination decreases depth by one.
The initial depth is $q+1$, proving the asserted depth bound.
The Mellin argument in Theorem~\ref{mx:poles} proves the continuation
and the entire assertion.  Equality on the original half-plane then
implies equality meromorphically everywhere and permits every fixed
$C$ derivative.
\end{proof}

\begin{lemma}[Eliminating a nonpositive trailing order]
\label{mx:eliminate}
Any multiple polylogarithm whose first order is $C+v$ with
$v\in\mathbb Z$ and whose remaining orders are integers reduces
finitely to terms of the form \eqref{mx:positiveoutput}, with depth
no greater than its initial depth.  Coincident cumulative colors are
handled directly, without division by their difference.
\end{lemma}

\begin{proof}
Choose any nonpositive trailing order $-M$ at position $j\geq2$.
For fixed neighboring indices put
$u=n_{j-1}$ and $l=n_{j+1}$, with $n_{k+1}=0$ when $j=k$ is last.
The sum to eliminate is
\[
 \sum_{l<n<u}\beta^n n^M,\qquad \beta=z_j.
\]
For $\beta\ne1$ define the rational functions and polynomial
\begin{equation}\label{mx:Q}
 S_r(\beta)=\left(\beta\frac d{d\beta}\right)^r
                         \frac1{1-\beta},\qquad
 Q_M(x;\beta)=\sum_{r=0}^M\binom Mr x^{M-r}S_r(\beta).
\end{equation}
The geometric-series identity, or its polynomial continuation in
$\beta$, gives
\begin{equation}\label{mx:geomelim}
 \sum_{l<n<u}\beta^n n^M
   =\beta^{l+1}Q_M(l+1;\beta)-\beta^uQ_M(u;\beta).
\end{equation}
For example, it follows by subtracting the two tails
$\sum_{n\geq v}\beta^n n^M=\beta^vQ_M(v;\beta)$ when
$|\beta|<1$, and then using equality of rational functions.  At
$\beta=1$ the formula used by the algorithm is instead
\begin{equation}\label{mx:faulhaber}
 \sum_{l<n<u}n^M
   =\frac{B_{M+1}(u)-B_{M+1}(l+1)}{M+1},
\end{equation}
where $B_1(x)=x-1/2$.  Both identities hold even when $u=l+1$;
their two boundary terms then cancel.

Expand the boundary polynomials in $u$ or $l$.  A power $u^a$ lowers
the preceding order by $a$, and the factor $\beta^u$ multiplies the
preceding color by $\beta$.  Similarly, a power $l^a$ lowers the next
order and $\beta^l$ multiplies its color by $\beta$.  At the last
position the lower boundary is a scalar.  After removing position
$j$, every term is a polylogarithm of depth one less, with first order
still of the form $C+$ an integer.  Some adjacent trailing orders
may now be nonpositive; repeat the same step.  The depth strictly
decreases at each step, so the recursion terminates after at most
$k-1$ eliminations along any branch.  Every trailing order in the
final terms is positive.  Equation~\eqref{mx:faulhaber} is the exact
coincident-color branch; it never requires taking numerically unstable
limits of \eqref{mx:geomelim}.
\end{proof}

\begin{remark}[What the reduction preserves]
Only the first polylogarithm order varies with $C$.  The theorem does
not authorize differentiating the finite expression with respect to
the discrete inner indices $s_j$.  Boundary elimination replaces a
pair of adjacent colors by their product; hence its cumulative colors
are obtained by deleting entries from
$(\alpha,b_1,\ldots,b_q)$.  In particular, if every original $X_j$
is nontrivial, it cannot introduce a cumulative color equal to one.
This observation also explains why equal neighboring colors are
harmless while a genuinely uncolored original factor can change the
outer spectral poles.
\end{remark}

The two unmixed sectors fit the same proof.  If $I_-=\varnothing$,
the shuffle gives a sum of depth-$d$ polylogarithms with first order
$C+k_1$.  If $I_+=\varnothing$, the coefficient identity
\eqref{mx:binomial} gives the preceding depth-one formula
\[
 Z(C)=\sum_{\alpha,h}c_{\alpha,h}
                    \sum_{r=0}^{h-1}e_{h,r}\mxLi_{C-r}(\alpha).
\]
Thus every pattern of integer inner orders has a finite reduction.

\subsection{Continuation, poles, and normal differentiation}

Positive uncolored factors introduce logarithms.  Therefore the
finite simple-pole list from the all-nonpositive rational sector
cannot simply be copied to the mixed sector.  The following local
rule gives all the poles and their principal coefficients.

\begin{theorem}[Complete principal-part prescription]
\label{mx:poles}
For the integer-order sum \eqref{mx:Z}, including the unmixed sectors,
put
\[
 F(x)=\prod_j\mxLi_{s_j}(X_je^{-p_jx}),\qquad
 M=\sum_{\substack{s_j\leq0\\X_j=1}}(1-s_j),\qquad
 \ell_+=\#\{j:s_j>0,\ X_j=1\}.
\]
Near $x=0^+$ there is a convergent logarithmic Laurent expansion
\begin{equation}\label{mx:locallaurent}
 F(x)=\sum_{v\geq-M}\sum_{b=0}^{\ell_+} a_{v,b}x^v(\log x)^b.
\end{equation}
The meromorphic continuation of $Z$ has possible poles only at
integers $m\leq M$.  Its complete principal part at $C=m$, with
$\varepsilon=C-m$, is
\begin{equation}\label{mx:principal}
 \operatorname{PP}_{C=m}Z(C)
 =\operatorname{PP}_{\varepsilon=0}\left\{
  \frac1{\Gamma(m+\varepsilon)}
    \sum_{b=0}^{\ell_+}\frac{(-1)^bb!\,a_{-m,b}}
                        {\varepsilon^{b+1}}\right\}.
\end{equation}
Consequently the pole order is at most $\ell_++1$ for positive $m$, and
at most $\ell_+$ for nonpositive $m$.  Missing coefficients and canceled
principal parts mean removable singularities.  If no $X_j=1$,
the function is entire.  Every normal derivative is obtained by
differentiating the meromorphic function and its displayed principal
parts.
\end{theorem}

\begin{proof}
For a nontrivial color the factor is analytic at zero, with expansion
\begin{equation}\label{mx:nontrivexp}
 \mxLi_s(Xe^{-px})
   =\sum_{k\geq0}\mxLi_{s-k}(X)\frac{(-px)^k}{k!}.
\end{equation}
For $a\geq0$ the uncolored nonpositive factor has expansion
\begin{equation}\label{mx:negativeexp}
 \mxLi_{-a}(e^{-px})
   =\frac{a!}{(px)^{a+1}}
      +\sum_{k\geq0}\zeta(-a-k)\frac{(-px)^k}{k!}.
\end{equation}
For $s\geq1$ the uncolored positive factor has the integer-order
limit of Jonqui\`ere's expansion~\cite[Eq.~25.12.12]{DLMFpolylog}:
\begin{equation}\label{mx:positiveexp}
 \mxLi_s(e^{-px})
  =\frac{(-px)^{s-1}}{(s-1)!}
       \{H_{s-1}-\log p-\log x\}
    +\sum_{\substack{k\geq0\\k\ne s-1}}
       \zeta(s-k)\frac{(-px)^k}{k!}.
\end{equation}
These convergent local expansions prove \eqref{mx:locallaurent} and
provide every needed coefficient by finite multiplication.

In the initial right half-plane, termwise integration gives
\begin{equation}\label{mx:mellin}
 Z(C)=\frac1{\Gamma(C)}\int_0^\infty x^{C-1}F(x)\,dx.
\end{equation}
The integral over any $[\delta,\infty)$ is entire before multiplication
by $1/\Gamma(C)$, since $F$ decreases exponentially at infinity.
Subtract successively more terms of \eqref{mx:locallaurent} near zero.
The remainder then defines a holomorphic function on successively
larger left half-planes, and the subtracted monomials contribute
\[
 \int_0^1x^{C+v-1}(\log x)^b\,dx
       =\frac{(-1)^bb!}{(C+v)^{b+1}}.
\]
Only $v=-m$ can contribute to the principal part at $C=m$.
The reciprocal Gamma function is nonzero at positive integers and
has a simple zero at each nonpositive integer.  This proves
\eqref{mx:principal} and the pole-order assertions.  If there are no
uncolored factors, then $M=\ell_+=0$; all potential Mellin poles cancel
against those Gamma zeros.  Local normal convergence of the subtracted
integrals on compact sets also justifies every fixed normal
derivative.

To verify continuation of each nested term separately, its coefficient
kernel, with the first denominator omitted, is
\[
 \sum_{n_1>\cdots>n_k\geq1}
 \frac{(z_1t)^{n_1}z_2^{n_2}\cdots z_k^{n_k}}
      {n_2^{t_2}\cdots n_k^{t_k}}
 =\frac{z_1t}{1-z_1t}
     \mxLi_{t_2,\ldots,t_k}(z_1z_2t,z_3,\ldots,z_k).
\]
For depth one the kernel is just $z_1t/(1-z_1t)$.  The positive
trailing indices make the right side a rational function times an
iterated integral.  Induction in its word length gives a convergent
finite-logarithmic local expansion at $t=1$.  Applying the same Mellin
subtractions with parameter $C+v$ proves its meromorphic continuation.
If none of its cumulative colors is one, its kernel is analytic at
$t=1$, and the continuation is entire.  As noted after
Lemma~\ref{mx:eliminate}, all these cumulative colors come from the
original root sets.
\end{proof}

Formula~\eqref{mx:principal} is fully finite at every specified pole.
For clarity, the Gamma factors needed there can be generated by
\begin{align}\label{mx:gammajets}
 \frac1{\Gamma(m+\varepsilon)}
 &=\frac1{(m-1)!}\exp\left\{
   (\gamma-H_{m-1})\varepsilon+
   \sum_{k\geq2}\frac{(-1)^{k+1}}k
       (\zeta(k)-H_{m-1}^{(k)})\varepsilon^k\right\},\quad m\geq1,
 \notag\\
 \frac1{\Gamma(-N+\varepsilon)}
 &=(-1)^NN!\varepsilon\exp\left\{
   (\gamma-H_N)\varepsilon+
   \sum_{k\geq2}\frac{(-1)^{k+1}\zeta(k)-H_N^{(k)}}k
                                        \varepsilon^k\right\}.
\end{align}
Only finitely many coefficients are required.  This controls all
principal coefficients through ordinary zeta values, generalized
harmonic numbers, integer-order polylogarithms, and logarithms of
the weights.  It does not assert that the holomorphic Laurent
coefficients reduce to depth one.

Inside the unit polydisk the functions are already entire in $C$.
When all boundary colors are nontrivial, their continuation agrees
locally uniformly in $C$ with the total-degree Abel limit
$X_j\mapsto\rho^{p_j}X_j$, $\rho\uparrow1$.  This follows from
the same Taylor subtractions, uniformly near those fixed colors.
If uncolored factors occur, an Abel limit outside the convergence
half-plane need not exist.  In that case the prescribed value is the
meromorphic continuation of the fixed-color sum, and a pole requires
a Laurent coefficient or other explicitly stated restriction.

\subsection{Explicit mixed and weighted identities}

For compactness write
\[
 Z_{s_1,\ldots,s_d}(C;z)
  =\sum_{m_1,\ldots,m_d\geq1}
      \frac{z^{m_1+\cdots+m_d}}
       {m_1^{s_1}\cdots m_d^{s_d}(m_1+\cdots+m_d)^C}.
\]
The following identities hold first for $|z|<1$, and thereafter with
the fixed-color continuation just specified.  From
$\mxLi_1(zt)^2=2\mxLi_{1,1}(zt,1)$ and
$\mxLi_0(zt)=-1+(1-zt)^{-1}$, one obtains
\begin{equation}\label{mx:110}
 Z_{1,1,0}(C;z)=2\mxLi_{C,1,1}(z,1,1).
\end{equation}
The repeated pole
\[
 \mxLi_{-1}(zt)=(1-zt)^{-2}-(1-zt)^{-1}
\]
gives, after the coincident-color elimination,
\begin{align}\label{mx:11minus1}
 Z_{1,1,-1}(C;z)={}&2\mxLi_{C-1,1,1}(z,1,1)
                   -2\mxLi_{C-1,1}(z,1)
                   +2\mxLi_{C,1}(z,1)\notag\\
                 &+2\mxLi_{C-1}(z)-2\mxLi_C(z).
\end{align}
Indeed the intermediate term is
$2\{\mxLi_{C-1,1,1}-\mxLi_{C,0,1}\}$, and direct summation of
the middle index proves
\[
 \mxLi_{a,0,1}(z,1,1)
 =\mxLi_{a-1,1}(z,1)-\mxLi_{a,1}(z,1)
      -\mxLi_{a-1}(z)+\mxLi_a(z).
\]
Every normal derivative of \eqref{mx:11minus1}, including at $C=0$
when $z\ne1$, follows by replacing each polylogarithm with its
corresponding first-order derivative.  This gives the whole normal
jet, rather than an isolated negative-integer value identity.

A four-variable example has depth at most three, as predicted:
\begin{align}\label{mx:1100}
 Z_{1,1,0,0}(C;z)={}&2\mxLi_{C-1,1,1}(z,1,1)
                    -2\mxLi_{C,1,1}(z,1,1)\notag\\
                  &-2\mxLi_{C-1,1}(z,1)
                    +2\mxLi_{C,1}(z,1)
                    +2\mxLi_{C-1}(z)-2\mxLi_C(z).
\end{align}
Here the rational factor is
$(zt)^2/(1-zt)^2=1-2/(1-zt)+(1-zt)^{-2}$.

An unequal-weight example, with both root filtering and a color
coincidence, is
\begin{equation}\label{mx:weightedexample}
 W(C)=\sum_{m,n,k\geq1}
       \frac{4^{-m}3^{-n}9^{-k}}{m^2n(2m+n+2k)^C}.
\end{equation}
It is entire in $C$.  Since
\[
 \mxLi_2(t^2/4)=2\{\mxLi_2(t/2)+\mxLi_2(-t/2)\},\qquad
 \frac{t^2/9}{1-t^2/9}
   =-1+\frac1{2(1-t/3)}+\frac1{2(1+t/3)},
\]
the three shuffles of the words $0a$ and $1/3$ give
\begin{align}\label{mx:weightedanswer}
 W(C)=\sum_{\substack{a\in\{1/2,-1/2\}\\
                     \delta\in\{1/3,-1/3\}}}
 \biggl\{&\mxLi_{C,1,2}(\delta,1/(3\delta),3a)
       +\mxLi_{C,2,1}(\delta,1/(3\delta),3a)\notag\\
       &+\mxLi_{C,2,1}(\delta,a/\delta,1/(3a))\biggr\}.
\end{align}
All cumulative colors in these twelve terms have modulus less than
one, so no continuation convention is needed to differentiate this
identity at any $C$.

The uncolored case of \eqref{mx:110} illustrates why the pole theorem
is necessary.  Near zero its Mellin kernel is
\[
 \mxLi_0(e^{-x})\mxLi_1(e^{-x})^2
 =x^{-1}\log^2x-\tfrac12\log^2x-\log x
   +x\{\tfrac1{12}\log^2x+\tfrac7{12}\log x+\tfrac14\}
   +O(x^2(1+|\log x|^2)).
\]
Consequently
\begin{align}\label{mx:examplepoles}
 \operatorname{PP}_{C=1} Z_{1,1,0}(C;1)
  &=\frac2{(C-1)^3}+\frac{2\gamma}{(C-1)^2}
                         +\frac{\gamma^2-\zeta(2)}{C-1},\notag\\
 \operatorname{PP}_{C=0} Z_{1,1,0}(C;1)
  &=-\frac1{C^2}+\frac{1-\gamma}{C},\notag\\
 \operatorname{PP}_{C=-1} Z_{1,1,0}(C;1)
  &=-\frac1{6(C+1)^2}+\frac{3/4-\gamma/6}{C+1}.
\end{align}
In fact there are infinitely many nonpositive poles: the coefficient
of $\log^2x$ in this kernel is exactly
$\mxLi_0(e^{-x})$, whose positive odd powers have nonzero Bernoulli
coefficients.  This observation concerns the mixed extension; it is
not a correction of the earlier theorem restricted to rational
generating functions.

\subsection{A compatible primitive hierarchy and exact checks}

With $r_n$ from \eqref{mx:factor}, introduce an additional variable
$z$ by $Z(C;z)=\sum_{n\geq1}r_nz^nn^{-C}$.  Equivalently this
replaces $X_j$ by $X_jz^{p_j}$.  In every term of the finite reduction
only the first polylogarithm color acquires the factor $z$.  Thus
\begin{equation}\label{mx:primitives}
 \int_0^z\partial_C^N Z(C;t)\frac{dt}{t}
        =\partial_C^N Z(C+1;z),\qquad
 z\frac d{dz}\partial_C^N Z(C;z)
        =\partial_C^N Z(C-1;z).
\end{equation}
The primitive is fixed by its zero value at $z=0$.  These identities
follow by local normal convergence in the disk and continue along
any common path avoiding the coefficient-function singularities.

The accompanying program performs 624 exact rational checks with
SymPy; its source is \texttt{verify\_mixed\_reduction.py}.
There are 484 coefficient
checks: for eleven independent fixtures, each of the first 22
coefficients is computed directly from the original product and
compared both with \eqref{mx:mainformula} and with the fully eliminated
positive-trailing-order expression.  The fixtures include four
original variables, three positive factors, unequal weights, repeated
rational poles, coincident and opposite unit colors, and nonpositive
intermediate trailing indices.  The remaining 140 checks verify the
finite boundary summations for degrees zero through six, including
empty intervals and $\beta=1$.  No floating-point tolerance is used.
The coefficient checks establish finite algebraic diagnostics;
meromorphic continuation and differentiation are supplied by the
proofs, not inferred from these tests.  The exact outputs and fixture
descriptions are recorded in \texttt{mixed\_reduction\_checks.json}.

The remaining question is arithmetic simplification of the resulting
positive-depth normal jets.  The depth bound is a guaranteed upper
bound, not an independence statement or a proof of minimality.
At roots of unity, one can now ask which symmetries or distribution
projections cancel these residual jets and leave only depth-one
Hurwitz-zeta, Gamma, and generalized-Stieltjes data.  Such a
cancellation problem is distinct from the finite mixed-sign closure
proved here.
